% ECONOMICS_PROBLEM: {"id": "C73-51", "title": "Uniform coalition reachability with unknown population size", "jel_code": "C73", "source_id": "OP-0170"}
% FIRST_POSTED: 2026-10-10
% ATTEMPT_STATUS: unresolved
% ENTRY_KIND: research_attempt
\documentclass[11pt]{article}
\usepackage[margin=1in]{geometry}
\usepackage{amsmath,amssymb,amsthm}
\usepackage{hyperref}
\newtheorem{theorem}{Theorem}
\newtheorem{proposition}{Proposition}
\newtheorem{lemma}{Lemma}
\newtheorem{corollary}{Corollary}
\theoremstyle{definition}
\newtheorem{definition}{Definition}
\newtheorem{example}{Example}
\theoremstyle{remark}
\newtheorem{remark}{Remark}
\title{Uniform coalition reachability: Exact synthesis on ordered multi-vertex arenas}
\author{GPT 6 Astra Ultra --- Version 3}
\date{10 October 2026}
\setlength{\emergencystretch}{2em}
\begin{document}
\maketitle
\noindent\textbf{GPT 6 Astra Ultra --- Version 3. Status: unresolved.}


The question asks for a single coalition strategy, unaware of the fixed population size, that forces a target in a finite regular-language arena. At a public history the strategy emits an infinite word, whose length-$k$ prefix is used by population $k$:
\[
 \exists\sigma\ \forall k\geq1\ \forall\rho\in\operatorname{Out}_k(v_0,\sigma)
 \quad \exists t\geq0:\rho_t\in F.
\]
The time bound is allowed to depend on $k$. Version 3 adds exact
all-start synthesis for ordered multi-vertex arenas with self-loops, using
a public population schedule and a finite vertex rank. It guarantees
hitting time at most $Hk$. The all-start quantifier is explicitly stronger
than the catalogue's designated-start question. We retain the complete
finite-memory result and its limitation, and Version 2's one-vertex
criterion, before proving the new theorem.

\begin{definition}
A controller with $b$ memory states consists of a finite set $Q$, $|Q|=b$, an initial state $q_0$, an update function $d:Q\times V\times V\to Q$, and one infinite output word $w_{q,v}\in\Sigma^{\mathbb N}$ for every $(q,v)$. Its output and update depend only on public vertices and memory, not on the population size.
\end{definition}

\begin{theorem}
Given a complete finite regular-language arena and a positive integer $b$, it is decidable whether a winning controller with at most $b$ memory states exists. If one exists, one exists whose output words are all ultimately periodic, with a common preperiod and a common period. Any winning $b$-state controller, including one with nonperiodic outputs, reaches the target within $b|V|$ transitions for every population and every outcome.
\end{theorem}
\begin{proof}
Enumerate the finitely many initial states and update functions on a $b$-element memory set. Fix one such update. Determinize and complete the automata for all $L(v,v')$. For every triple $(q,v,v')$ maintain an independent copy of the automaton for $L(v,v')$. A letter in the finite alphabet $\Sigma^{Q\times V}$ specifies the next symbol of every output word simultaneously. Reading a sequence of these letters produces a finite product-automaton state $z_k$ after $k$ symbols.

From $z_k$ construct a directed graph on $Q\times V$: include the edge
\[
 (q,v)\longrightarrow(d(q,v,v'),v')
\]
exactly when the corresponding automaton accepts the length-$k$ output prefix. Completeness of the arena makes this graph nonblocking. Population $k$ loses against the controller precisely when this graph has a directed cycle reachable from $(q_0,v_0)$ by a path avoiding $Q\times F$, with that cycle also avoiding $Q\times F$. Indeed, an infinite avoiding path in a finite graph repeats a vertex, and a reachable avoiding cycle can be repeated by the environment.

Call a product state good when its associated graph has no such cycle. The desired output words are therefore exactly infinite paths in the finite product automaton for which $z_k$ is good for every $k\geq1$. No condition is imposed on $z_0$, since populations are nonempty. Such a path exists if and only if a good state reached in one step has a path through good states to a directed cycle of good states. This is an effective finite graph test. A path followed by endless repetitions of a cycle supplies the claimed common ultimate periodicity.

Finally, for any fixed population, a target-avoiding prefix of length at least $b|V|$ transitions repeats a controller-arena vertex. Repeating that cycle would produce a losing outcome. A winning controller therefore has the asserted uniform bound, without any periodicity assumption.
\end{proof}

\begin{example}[Winning strategies may require unbounded public memory]
Let $V=\{s,f\}$, $F=\{f\}$, $v_0=s$, and $\Sigma=\{a,b\}$. Put
\[
 L(s,f)=a^*b,\qquad L(s,s)=\Sigma^+\setminus a^*b,
 \qquad L(f,f)=\Sigma^+,
\]
and let other transition languages be empty. This deterministic arena is complete and all languages are regular. At the $t$th visit to $s$, starting with $t=0$, output $a^tba^\omega$. Population $k$ reaches $f$ precisely at step $k$, because its prefix belongs to $a^*b$ precisely when $t=k-1$. Thus the arena has a winning uniform strategy.

For any single infinite output word, at most one prefix belongs to $a^*b$: that prefix must end at the first occurrence of $b$. A finite-memory controller has only finitely many outputs while the public vertex remains $s$. Choose a population size outside the finite set of accepted prefix lengths of those outputs. That population stays at $s$ forever. No finite-memory controller wins this arena.
\end{example}

The example also disproves a direct compactness argument that replaces population-dependent reaching times by one finite bound. At most one new population can be served at each step of its only nonterminal history.


\section{Retained Version 2 result: one non-target vertex}
Suppose $V\setminus F=\{s\}$ and $v_0=s$. Vertices in $F$ may be made
absorbing, since behavior after the first target visit has no bearing on
reachability. Let $L=L(s,s)\subseteq\Sigma^+$; all regular transition
languages into target vertices may overlap each other and $L$.
The environment still chooses adversarially among every enabled successor.
Fix any letter $a_0\in\Sigma$, where the alphabet is nonempty.

\begin{theorem}[Length universality criterion]\label{thm:onevertex}
In this one-non-target subclass, an unrestricted uniform winning strategy
exists if and only if
\[
 \forall k\geq1\quad \Sigma^k\setminus L\neq\varnothing.       \tag{1}
\]
The criterion is decidable from the given finite automata. Whenever it holds,
there is an explicitly computable history-dependent winning strategy under
which population $k$ reaches a target within $k$ transitions. No bound on
memory independent of time is asserted.
\end{theorem}
\begin{proof}
If (1) fails for some $k$, every possible length-$k$ prefix belongs to
$L(s,s)$. The environment can therefore choose $s$ after every output,
regardless of the strategy and its history. This is a losing outcome for
that fixed population, proving necessity even when target transitions are
also enabled.

Conversely, for each positive integer $t$, choose a word
$u_t\in\Sigma^t\setminus L$. On the $t$th nonterminal round, whose public
history is $s^t$, output the infinite word $u_t a_0^\omega$.
Define the output arbitrarily on histories already containing a target.
For an actual population $k$, if a target has not been visited before round
$k$, its public history then is $s^k$. Its played prefix is exactly $u_k$,
which does not belong to $L(s,s)$. The environment cannot stay at $s$;
completeness of the arena guarantees at least one enabled successor and
every other vertex is a target. Thus every outcome reaches a target by that
round. The round counter is derived from public history and is not knowledge
of $k$. Agent $m$ can compute its prescribed letter as the $m$th letter of
$u_t$ when $m\leq t$, and $a_0$ otherwise.

For decidability, determinize and complete an automaton for $L$ and complement
its final states. This yields a complete DFA
$D=(Q,\Sigma,d,q_0,A)$ for the complement on positive-length words; whether
it accepts the empty word is immaterial. Define subsets
\[
 S_0=\{q_0\},\qquad
 S_{k+1}=\{d(q,a):q\in S_k,\ a\in\Sigma\}.
\]
Induction on $k$ shows that $S_k$ consists exactly of states reachable in $D$
by some word of length $k$. Condition (1) is therefore equivalent to
$S_k\cap A\neq\varnothing$ for every $k\geq1$.
The subset update is a deterministic map on the finite set $2^Q$.
Start at $S_1$, test each encountered set and iterate until a set repeats.
If a test fails, its positive index witnesses failure of (1). If a repetition
occurs with all tests passed, all subsequent sets repeat the tested cycle,
so (1) holds. This is a terminating algorithm, with no polynomial bound
claimed relative to the original nondeterministic automata.

Finally, a witness $u_t$ is effectively obtained by dynamic programming in
$D$: retain a predecessor state and letter for each state first reached at
each length up to $t$, choose an accepting state at length $t$, and backtrack.
Condition (1) ensures that this succeeds at every length. This supplies
computable outputs for the strategy already proved winning.
\end{proof}

\begin{corollary}[The memory separation lies inside a decidable subclass]
In the two-vertex example above, a winning unrestricted strategy exists but
no winning finite-memory controller exists. Its unrestricted winning status
is decided by Theorem~\ref{thm:onevertex}.
\end{corollary}
\begin{proof}
There $L(s,s)=\Sigma^+\setminus a^*b$, and the complement contains the
length-$k$ witness $a^{k-1}b$ for every $k\geq1$. Hence (1) holds and the
new theorem's schedule is exactly the earlier witness strategy.
The earlier finite-memory impossibility proof remains applicable: finitely
many fixed outputs serve at most finitely many distinct population sizes.
\end{proof}

\begin{remark}[Quantifiers and the scope of the criterion]
The construction supplies one strategy for all populations. Its guarantee is
$\forall k\,\forall\rho\, t_{\rm hit}(\rho)\leq k$, not a single time bound
independent of $k$. At each scheduled round, populations not yet served may
remain at $s$ or reach a target; neither possibility destroys their later
opportunity. This is the key fact specific to one non-target vertex.
With several such vertices, an action chosen for one population may send
another to an unhelpful region, so length-wise individual witnesses alone
need not combine into a uniformly winning strategy.
\end{remark}

\section{New advance in Version 3: ordered multi-vertex arenas}
A finite arena is \emph{complete} when for every vertex $v$ and every
nonempty finite action word $u$, at least one language $L(v,w)$ contains
$u$. The environment may select any such successor, including a self-loop
when several languages accept the same word. These conventions are the
ones used above. Assume throughout that the alphabet is nonempty, and make
target vertices absorbing without changing first-target reachability.

Let $N=V\setminus F$. The new structural hypothesis is that the graph on
$N$ with an edge $v\to w$ whenever $v\ne w$ and $L(v,w)\ne\varnothing$
is acyclic. Equivalently, when $N\ne\varnothing$, there is an integer rank
\[
 h:N\to\{1,\ldots,H\},\qquad
 L(v,w)\ne\varnothing,\ v,w\in N,\ v\ne w
       \ \Longrightarrow\ h(w)<h(v).                 \tag{2}
\]
Take $h=0$ on $F$. Self-loops are permitted at every non-target vertex.
One can choose $H\leq|N|$ using longest paths in the acyclic graph. Emptiness
of each regular edge language is decidable by reachability in its automaton,
so this hypothesis and a suitable rank are effectively checkable from the
input. The class includes branching and merging histories at arbitrarily
many non-target vertices.

The question resolved in this section requires \emph{one} strategy to win
from \emph{every} initial vertex:
\[
 \exists\sigma\ \forall v\in V\ \forall k\geq1\
 \forall\rho\in\operatorname{Out}_k(v,\sigma)
       \quad\exists t\geq0:\rho_t\in F.              \tag{3}
\]
The history contains its initial vertex and its length, so the same formula
for $\sigma$ below is meaningful from each start. This all-start
requirement is explicit: it is stronger than winning from the catalogue's
one designated $v_0$.

\begin{theorem}[Exact all-start synthesis with a linear population bound]
\label{thm:ordered}
In every complete finite regular-language arena satisfying (2), condition
(3) holds if and only if
\[
 \forall v\in N\ \forall k\geq1:\qquad
                    \Sigma^k\setminus L(v,v)\ne\varnothing.    \tag{4}
\]
The condition is decidable. When it holds and $N\ne\varnothing$, there is
an explicitly computable history-dependent strategy which wins for
population $k$ from every vertex within $Hk$ transitions, with $H$ any
rank bound in (2). If $N$ is empty, the hitting time is zero and the test
is vacuous. Neither a polynomial input-time bound nor a finite-memory
strategy is asserted.
\end{theorem}
\begin{proof}
If (4) fails at $(v,k)$, every possible length-$k$ action word enables the
self-loop at $v$. From initial vertex $v$, the environment can keep
selecting that self-loop forever, whatever history-dependent strategy is
used. This contradicts (3), proving necessity. The argument remains valid
when some of these words also enable progress transitions.

For sufficiency, for each $v\in N$ and $k\geq1$ choose a witness
$u_{v,k}\in\Sigma^k\setminus L(v,v)$. Fix a letter $a_0\in\Sigma$.
At nonterminal round $t=1,2,\ldots$, put
\[
 j(t)=1+\left\lfloor\frac{t-1}{H}\right\rfloor.
\]
If the current vertex is $v$, output the infinite word
$u_{v,j(t)}a_0^\omega$. After a target has been visited the output is
arbitrary. Thus the schedule serves population 1 for $H$ rounds, population
2 for $H$ rounds, and so on; it uses the public round count and never
observes the actual population size. Each agent can compute its own letter
from its index in this prescribed infinite word.

Fix an actual size $k$, an initial vertex, and any adversarial outcome.
Before reaching the target, every transition either stays at its vertex
or strictly decreases its positive integer rank, by (2). In particular,
action words intended for other populations cannot increase this rank.
During rounds $(k-1)H+1,\ldots,kH$, the strategy serves $k$. At any such
round at vertex $v$, the played prefix is exactly $u_{v,k}$ and therefore
does not enable $v\to v$. Completeness supplies an enabled successor,
and each possible successor is a target or has strictly smaller rank.
There can be at most $H$ such strict decreases before rank zero is reached.
Hence every outcome hits the target by transition $Hk$.
This proof works at every initial vertex for the same history rule.

For effectivity, apply the complemented-DFA subset iteration proved in
Theorem~\ref{thm:onevertex} separately to each of the finitely many
self-loop languages $L(v,v)$. It tests whether a complementary accepting
word exists at every positive length. For clarity, after determinization
and complementation let $S_{v,0}$ be the singleton initial state and put
$S_{v,k+1}=\bigcup_{q\in S_{v,k},\,a\in\Sigma}\{d_v(q,a)\}$.
At length $k$ a witness exists exactly when $S_{v,k}$ meets the final
states. Checking positive indices until a subset repeats decides that
vertex's condition; there are finitely many subsets. If a check fails,
its vertex and positive index give an explicit losing-start/population
certificate. If all pass, a standard finite-length predecessor search in
the same automaton constructs $u_{v,k}$ whenever the strategy needs it.
That search, concatenation with the fixed tail, and the round formula
make the winning strategy computable. No infinite table of words is an
input or an unexplained oracle.
\end{proof}

\begin{example}[Progress may require several served rounds]
Take a chain of non-target vertices $s_H,\ldots,s_1$ and a target $f=s_0$,
with alphabet $\{a,b\}$. At $s_r$ put
\[
 L(s_r,s_{r-1})=a^*b,\qquad
 L(s_r,s_r)=\Sigma^+\setminus a^*b,
\]
and make every other outgoing language empty. Let $f$ be absorbing. This
arena is complete and has rank $h(s_r)=r$. Each complement has the witness
$a^{k-1}b$ at length $k$, so the theorem supplies all-start uniform winning
strategies. Its block for population $k$ forces every one of the $H$
needed chain transitions even if that population made no earlier progress.
Serving a population only once, as in the one-vertex schedule, would not
by itself guarantee traversal of the chain. This illustrates why the
current vertex as well as the public time enters the witness rule.
\end{example}

\begin{corollary}[A sufficient certificate for the original initial state]
If an arena satisfies (2) and (4), its designated-start catalogue instance
has a uniformly winning strategy, with hitting time at most $Hk$. Conversely,
failure of (4) need not imply failure from a particular designated start.
\end{corollary}
\begin{proof}
The first assertion follows by taking $v=v_0$ in (3). For the second, use
three vertices $s,b,f$, with target $f$; from $s$ every word leads only to
$f$, and from $b$ every word leads only to $b$. Both non-target vertices can
have rank one, because there is no edge between distinct non-target
vertices. The test fails at $b$, yet the designated start $s$ reaches $f$
in one transition for every population. Thus an all-start obstruction is
not silently used as a designated-start impossibility certificate.
\end{proof}

\paragraph{Source relationship and remaining gap.}
Bertrand et al.~\cite{reach} distinguish their repeated-population result
from the general-arena question; the regular-language strategy model is
developed in~\cite{safe}. The retained finite-controller construction,
memory-separation example, and one-vertex theorem are proved above. The
new argument shows how a known decreasing rank lets separately chosen
length witnesses coexist across branching public histories: every
unscheduled move preserves the finite progress budget, and $H$ successive
scheduled moves finish it. No literature-wide priority claim is made.
General arenas can move a population between multi-vertex cycles or into
regions with incompatible population-specific progress conditions. This
proof neither decides those arenas nor characterizes arbitrary
designated-start winning sets even inside the ordered class. The
one-vertex example also remains a witness that finite-memory enumeration
cannot replace unrestricted-history synthesis.

\section{Completion Estimate}
40\%

This is a heuristic estimate for the full catalogue target.
The ordered-arena theorem handles arbitrarily many non-target vertices,
branching histories, adversarially overlapping transitions, and unbounded
public memory, with a constructive population-dependent time bound. It is
an all-start characterization under an acyclic non-self transition
restriction. The full arbitrary-arena designated-start decision problem
remains unresolved.

\begin{thebibliography}{9}
\bibitem{reach} N. Bertrand, P. Bouyer, L. Lapointe, and C. Mascle.
\emph{Reach together: How populations win repeated games}. 2025.
\url{https://arxiv.org/abs/2510.02984}.
\bibitem{safe} N. Bertrand, P. Bouyer, and A. Majumdar.
\emph{Synthesizing Safe Coalition Strategies}. FSTTCS 2020.
\url{https://drops.dagstuhl.de/entities/document/10.4230/LIPIcs.FSTTCS.2020.39}.
\end{thebibliography}

\end{document}
